\section{Introduction}\label{sec:introduction}

The central failure mode in recursive multi-agent systems is not malfunction.
It is \emph{drift}: the gradual, often invisible separation of a spawned agent's
behavior from the intent of its originating principal, compounded across
generations of recursive spawning until the system's observable trajectory
bears little recognizable relation to any human-directed purpose.  Drift does
not require a faulty agent.  It emerges from the structure of accountability
chains themselves — specifically, from the absence of structural constraints that
prevent intermediate nodes from severing or obscuring the causal link between
a leaf agent's actions and the human principal who authorized its ancestry.

In uncontrolled multi-agent deployments, each spawning event creates a new
accountability boundary.  The parent agent — not the human principal — becomes
the immediate authority over the child.  If the parent is compromised,
misaligned, or simply operating under conditions of partial observability,
its delegation carries forward the full weight of its uncertainty.  The child
agent inherits not only a task but an accountability posture: a set of
implicit assumptions about objectives, constraints, and escalation paths that
may never have been explicitly reviewed by a human.  When that child spawns
its own descendants, the compounding is structural.  By the third or fourth
generation, the human principal may be factually excluded from the causal
chain — not by design, but by the accumulated weight of delegation chains
that no layer was structurally required to anchor.

The field has attempted to address drift through two families of
shallow-restriction mechanisms.  The first is \textbf{time-to-live} (TTL)
enforcement: spawn chains are permitted to extend only to a fixed depth or
wall-clock limit, after which further spawning is blocked.  The second is
\textbf{depth limits} with static capability caps: agents above a given
generation threshold lose access to certain tool classes, cannot write to
certain state partitions, or require additional confirmation before acting.
Both approaches are architecturally sound as circuit breakers but are
fundamentally insufficient as accountability mechanisms.  TTL enforcement
does not distinguish between a spawn chain that has correctly executed its
purpose and one that has diverged from it — both terminate at the same depth.
Depth limits treat capability restriction as a proxy for accountability, but
capability restriction does not prevent an agent from taking many low-capability
actions that collectively produce a high-capability outcome.  Neither approach
addresses the root structural condition: the absence of an immutable,
verifiable chain connecting every node's authorization state to a named human
principal.

The core structural insight of the \bxthree{} Recursive Spawning pillar is
that drift is made structurally impossible by one mechanism alone: an
\textbf{immutable parent pointer chain} that is enforced by the \Fact{} layer
of every node, such that no agent — including the node's own parent — can
modify, truncate, or obscure the chain of accountable ancestry linking any
node to its originating human principal.  In the \bxthree{} three-layer model
(\Purpose{}, \Bounds{}, \Fact{}), each layer has a distinct function in this
enforcement.  The \Purpose{} layer carries the human principal's intent and
the canonical authorization scope for the node.  The \Bounds{} layer constrains
the node's reasoning and action space to provisioned operating ranges.  The
\Fact{} layer is the enforcement substrate: it is where the immutable parent
pointer is written at spawn time, where all spawn authorizations are recorded
with cryptographic attribution, and where any attempt to reference a modified
or severed ancestry chain triggers a hard abort routed directly to the
upstream human principal via the Bailout Protocol.

We define drift formally as follows:

\begin{invariant}[Drift Condition]\label{invariant:drift}
A \bxthree{} node $n$ is in a state of \emph{drift} at time $t$ if and only if
there exists a causal chain of agent actions $\alpha_1, \alpha_2, \dots,
\alpha_k$ such that:
\begin{enumerate}
  \item $\alpha_1$ was authorized by a human principal $h$;
  \item each subsequent $\alpha_{i+1}$ was taken by a node spawned (directly or
    transitively) as a descendant of the node that took $\alpha_i$; and
  \item the \Fact{} layer of the node taking $\alpha_k$ does not contain a
    verifiable parent pointer chain traceable to $h$.
\end{enumerate}
\end{invariant}

Under this definition, drift is not a behavioral condition — it is a
structural one.  A node may exhibit drift without any individual agent acting
maliciously or incorrectly.  Drift is detected not by monitoring behavioral
outputs but by inspecting the integrity of the parent pointer chain maintained
in the \Fact{} layer.  The \bxthree{} Recursive Spawning pillar guarantees:

\begin{property}[Drift Impossibility]\label{prop:drift-impossibility}
In a \bxthree{} system in which the Recursive Spawning pillar is correctly
implemented, no node can exist in the state defined in
Invariant~\ref{invariant:drift}.
\end{property}

This guarantee holds because the immutable parent pointer chain is written
atomically at spawn time by the \Fact{} layer of the parent node, and this
write is non-repudiable: the parent cannot modify, revoke, or obscure the
pointer after the child has been instantiated.  The child's \Fact{} layer
validates its own ancestry on every decision cycle.  If the chain is broken
or untraceable to a named human principal, the node triggers the Bailout
Protocol immediately, without discretionary routing through the parent.

The Recursive Spawning pillar operates in conjunction with the other four
\bxthree{} pillars.  Loop Isolation prevents circular spawn graphs from
creating ambiguous ancestry states.  Spatial Firewall ensures that spawned
agents execute in execution domains distinct from their parents', so that
accountability boundaries map cleanly to resource boundaries.  Sandbox Gate
governs the communication channels between parent and child domains,
enforcing that inter-domain messages pass through authorized channels with
attested payloads.  Bailout Protocol is the terminal enforcement mechanism:
when drift is detected or when any exception condition cannot be resolved
within the node's authorized scope, the exception state propagates directly
to the named human principal with full forensic attribution at every hop.

The remainder of this paper is organized as follows.  Section~
\ref{sec:background} situates the Recursive Spawning pillar within the
broader literature on multi-agent accountability and recursive system safety.
Section~\ref{sec:three-layer} formally introduces the \bxthree{} three-layer
model as the architectural substrate.  Section~\ref{sec:pillars} provides the
full formal specification of all five pillars, with detailed treatment of the
immutable parent pointer chain mechanism, the spawn authorization protocol,
and the ancestry verification routine executed by the \Fact{} layer on each
decision cycle.  Section~\ref{sec:correctness} establishes the correctness
properties guaranteed by the pillar, including the drift impossibility theorem
and the characterization of the pillar's behavior under adversarial spawn
conditions.  Section~\ref{sec:related} evaluates the Recursive Spawning
approach against existing spawning governance models including TTL enforcement,
capability-tiered depth limits, and dynamic delegation frameworks.
Section~\ref{sec:limitations} addresses structural limitations and the
directed research agenda.  Section~\ref{sec:conclusion} concludes.
